% ============================================================
%  第十一章  安装部署与运行
% ============================================================
\section{安装部署与运行}

\subsection{环境准备}

推荐使用 Conda 创建独立环境，避免与系统 Python 包冲突。本项目统一使用名为 \code{drone} 的环境：

\begin{lstlisting}[style=sh, caption={创建并激活虚拟环境}]
conda create -n drone python=3.11 -y
conda activate drone
\end{lstlisting}

选择 Python 3.11 的原因：PyTorch、OpenCV、ultralytics 等依赖在该版本上的预编译轮子最完整；过新的版本（如 3.13/3.14）常出现轮子缺失、需本地编译的问题。

\subsection{依赖安装}

\subsubsection{安装 PyTorch}

先安装与硬件匹配的 PyTorch，再安装其余依赖，可避免 pip 自动拉取不匹配的版本：

\begin{lstlisting}[style=sh, caption={PyTorch 安装（按硬件三选一）}]
# CPU 版本
pip install torch torchvision --index-url https://download.pytorch.org/whl/cpu

# NVIDIA GPU（CUDA 12.4）
pip install torch torchvision --index-url https://download.pytorch.org/whl/cu124

# NVIDIA GPU（CUDA 11.8）
pip install torch torchvision --index-url https://download.pytorch.org/whl/cu118
\end{lstlisting}

\subsubsection{安装其余依赖}

\begin{lstlisting}[style=sh, caption={安装项目依赖}]
cd D:\CodeWorks\drone
pip install -r requirements.txt
\end{lstlisting}

\file{requirements.txt} 的内容为：flask、requests、pywebview、ultralytics、opencv-python、numpy、tqdm、matplotlib。

\subsubsection{版本验证}

安装完成后执行自检：

\begin{lstlisting}[style=sh, caption={环境自检}]
python -c "import flask, cv2, torch, numpy, ultralytics, webview; \
print(flask.__version__, cv2.__version__, torch.__version__)"
\end{lstlisting}

本次文档编写环境实测的版本如表~\ref{tab:versions} 所示。

\begin{table}[H]
\centering
\caption{实测环境版本}
\label{tab:versions}
\small
\begin{tabularx}{\linewidth}{@{}>{\heiti}p{3.4cm}>{\ttfamily}p{4.2cm}>{\raggedright\arraybackslash}X@{}}
\toprule
组件 & 实测版本 & 说明 \\
\midrule
操作系统 & Windows 11（x64） & 亦支持 Windows 10 \\
Python & 3.11.17 & Conda 环境 \code{drone} \\
Flask & 3.1.3 & Web 框架 \\
PyTorch & 2.14.1+cpu & CPU 轮子，GPU 环境为 +cu124 \\
torchvision & 0.29.1+cpu & 与 PyTorch 配套 \\
OpenCV & 5.0.0 & 视频解码与图像处理 \\
NumPy & 2.4.6 & 数组与文件读写 \\
ultralytics & 8.4.191 & YOLO 推理与跟踪 \\
pywebview & 5.x & 桌面窗口（依赖 WebView2） \\
\bottomrule
\end{tabularx}
\end{table}

\begin{notebox}[WebView2 运行时]
Windows 11 已内置 WebView2 运行时；Windows 10 若缺少，需安装“Microsoft Edge WebView2 Runtime”。若不希望安装该运行时，可使用 \code{--headless} 模式，通过浏览器访问。
\end{notebox}

\subsection{模型与数据准备}

\begin{enumerate}
  \item \term{YOLO 权重}：放置在 \file{models/yolo/}，默认读取 \file{detect.pt}。首次使用 ultralytics 时若指定 \code{yolov8n.pt} 等官方名称，会自动从网络下载。
  \item \term{预测权重}：放置在 \file{models/kinflow/}（或训练产物 \file{train/kinflow_predictor/train_result/best/}）。按“\code{best_s2} $\rightarrow$ \code{best_s1}”的优先级选择；缺失时自动回退线性外推。
  \item \term{轨迹数据}：\file{data/fact/} 下的 \file{visible.npz}、\file{infrared.npz}、\file{radar.npz}。首次运行若为空，界面显示空场景，不影响启动。
\end{enumerate}

\subsection{启动与访问}

\begin{lstlisting}[style=sh, caption={启动方式}]
python main.py               # 桌面窗口，同时启动 5000 与 5001
python main.py --headless    # 纯 HTTP，浏览器访问
python main.py recon         # 仅轨迹重建与分析
python main.py recog         # 仅无人机检测
\end{lstlisting}

启动成功后控制台输出两个地址；浏览器访问 \code{http://127.0.0.1:5000} 与 \code{http://127.0.0.1:5001} 即可。桌面模式下，重建模块以 1280$\times$800 的窗口打开，检测模块不弹窗，通过导航栏“识别”入口跳转。

启动前程序会扫描并终止占用 5000/5001 的残留进程（Windows 使用 \code{netstat} + \code{taskkill}，类 Unix 使用 \code{lsof} + \code{kill}），因此重复启动不会因端口占用失败。

\subsection{部署建议}

\begin{table}[H]
\centering
\caption{不同使用场景的部署建议}
\label{tab:deploy}
\small
\begin{tabularx}{\linewidth}{@{}>{\heiti}p{3.4cm}>{\raggedright\arraybackslash}X@{}}
\toprule
场景 & 建议 \\
\midrule
单机演示 & 默认桌面模式，双击启动脚本即可；无需联网（AI 策略除外） \\
现场指挥 & \code{--headless} 模式，两个标签页分别打开 5000/5001；关闭屏幕保护 \\
GPU 加速 & 安装 CUDA 版 PyTorch，将 \code{detection.device} 设为 \code{cuda:0} \\
离线环境 & 预先下载模型权重并放入 \file{models/}；前端库已随包提供，无需 CDN \\
二次开发 & 以 \code{flask run} 方式启动并开启调试；或使用 \code{python -m} 方式分模块调试 \\
\bottomrule
\end{tabularx}
\end{table}

\subsection{常见问题排查}

\begin{table}[H]
\centering
\caption{常见问题与处理}
\label{tab:faq}
\small
\begin{tabularx}{\linewidth}{@{}>{\heiti}p{4.0cm}>{\raggedright\arraybackslash}X@{}}
\toprule
现象 & 原因与处理 \\
\midrule
端口被占用 & 程序会自动清理；若仍失败，手动执行 \code{netstat -ano | findstr :5000} 后终止进程 \\
页面空白 / 样式缺失 & 检查静态资源路由；浏览器强制刷新（Ctrl+F5） \\
无法打开桌面窗口 & 缺少 WebView2 运行时，改用 \code{--headless} 模式 \\
检测无结果 & 模型类别与目标不符；确认 \code{target_classes} 与模型类别一致 \\
预测点成直线 & 混合模型未加载，正在使用线性外推（见第~\ref{sec:prediction_future} 章） \\
AI 策略报错 & 检查 \code{config.json} 中 \code{ai} 段的地址、密钥与模型名；确认网络可达 \\
配置被改坏 & 删除 \file{config.json}，重启后自动从模板重建 \\
\bottomrule
\end{tabularx}
\end{table}
